\chapter{绪论}


\section{研究背景与意义}

近年来，随着网易云音乐、QQ音乐等中心化流媒体分发平台的快速发展，音乐作品的传播和消费形式也发生了显著变化。然而，中心化数字音乐产业内长期存在的若干结构性问题却并未得到有效解决，主要体现在版权确权可信度不足、授权链复杂、收益分配透明度低等方面。对于独立音乐人和小型团队而言，音乐作品的发行、授权管理、收益核验和粉丝运营依赖于平台审核，创作者对作品实际传播路径和收益分配的掌控权较小。与此同时，版权凭证也较难在平台间实现统一验证\cite{Torbensen2019,Arenal2024}，用户能访问音乐，但很难验证访问权限是否可追溯。

从技术角度看，区块链技术的发展为数字音乐版权管理提供了一种全新的可信记录机制。因其具有可追溯、可编程以及不可篡改的特性，区块链适合用于记录作品授权状态、交易过程和收益规则等对可信度要求较高的场景。与此同时，IPFS 技术提供了一种以内容寻址方式保存文件的机制，以内容哈希（CID）作为资源标识，为音轨、封面、歌词、元数据等大型多媒体资源提供了链下的去中心化存储方案，从而形成“链上记录可信状态、链下保存内容资源”的DApp分层架构。此外，类似 SIWE（Sign-In with Ethereum）的标准消息签名身份认证机制的成熟，使得用户能够以较低门槛实现 Web3 身份登录，从而让开发者能够将链上身份与链下会话管理结合起来，降低区块链系统的开发成本和使用门槛\cite{ipfsDocs,eip4361}。

完全去中心化的架构存在治理复杂度高、可扩展性差等问题\cite{Torbensen2019,Arenal2024,Chen2021,Rikken2019}，同时，高频业务数据上链也会导致区块链交易成本高昂、存取性能较低\cite{Wei2022BDMSSurvey,Liu2020GasExceptions}。本文设计并实现了 FreeFlow 这一基于区块链的去中心化音乐社区平台原型，旨在解决传统中心化流媒体平台版税分配不透明、版权凭证难验证等问题。在工程实现上，FreeFlow 放弃了完整社区自主治理和全量数据上链的全 Web3 解决方案，转而采用“链上高可信确权 + 链下高效服务”的 Web2.5 混合架构。系统以 ERC-1155 合约接口作为资产标准\cite{eip1155}，链上负责统一管理作品发布、访问凭证与版税规则，而评论、链上数据检索、会话管理等高频业务数据则保留在链下后端服务器，以权衡用户体验与运维成本。


\section{国内外研究现状}

\subsection{数字音乐发行系统研究现状}

现有的数字音乐分发平台，大多都采用中心化的架构，来对内容访问、权限控制以及收益分配进行统一的管理。其版税分配，通常选用\textit{pro-rata}模式，也就是建立起一个统一的收入池，并按照播放量等数据，在该池内进行收益分配\cite{Alaei2022ProRata,Moreau2024PaymentModels}。相关研究指出，在这种模式下，作品的录音、创作者、出版者还有权利份额等信息，是由多方共同参与的，使得权利链条变得很分散；并且，版税的分配又常常需要依靠多层中介机构，创作者难以了解作品实际的传播和收入状况\cite{Torbensen2019,Arenal2024}。研究还表明，元数据的不一致、许可结构的复杂性，以及版税发放效率的低下，已经成为了数字音乐产业当中最具代表性的结构性问题\cite{Torbensen2019,Arenal2024}。

国内方面，在区块链和版权管理方面的研究更多强调版权存证与司法协同，以提升司法处理效率\cite{chinaCopyrightChain2021,beijingCopyrightChain2022}。而对去中心化流媒体分发平台本身的研究较少。这也说明，在国内构建一套基于区块链的服务，往往需要兼顾功能实现和可取证性\cite{chinaCopyrightChain2021,beijingCopyrightChain2022}。

综上，现有数字音乐发行系统研究在版税分配、流媒体服务和版权管理方面已形成较多成果，但对“发布—授权—访问—互动—发现”的全流程一体化设计关注仍然有限，仍缺乏一个将创作者自主发行、链上访问控制、用户社区和音乐检索有机统一的平台\cite{Li2021,Niu2022,Shi2025}。

\subsection{音乐检索与资源排序研究现状}

音乐检索与资源排序是数字音乐平台中的重要模块。构建一套高效、可解释的算法，有助于用户从海量音乐资源中快速检索出贴近需求的结果。传统音乐检索算法通常包括基于标题、艺人、专辑、标签和描述文本的元数据检索，也包括像 Shazam、听歌识曲等基于音频内容特征的检索方法。在如今的大型中心化平台中，算法往往能够结合多元数据，并借助用户行为分析等多方面信息，为用户提供搜索服务。而对于 FreeFlow 这种原型性质的平台来说，用户行为数据比较有限，复杂的个性化模型和算法难以充分发挥作用。但是，基于可解释特征的检索排序方法，仍然具有较高的实用价值。

从研究的积累状况来看，在音乐信息检索领域，已经形成了一套相对成熟的数据集和评价指标体系。像 Million Song Dataset、FMA 还有 LFM-1b 等数据集，分别从大规模音频特征、开放音乐元数据以及听歌行为事件这些方面，给音乐检索、召回和推荐方面的研究提供了支持；而 MusDB18 等数据集，则更加适宜用于音频分离、水印检测以及音频处理等任务。这也就表明，“音乐平台系统研究”和“音乐算法评测”这两者，在数据和任务上存在着明显的差异：前者更加侧重于系统流程和功能的闭环，后者则更加侧重于检索、推荐或者音频分析算法的评价设计。因此，若系统包含独立的资源检索与排序模块，就有必要借鉴信息检索领域较成熟的评价方法，并结合相关性指标与离线评估判断模块效果\cite{msd2011,fma2017,lfm1b2016}。

结合 FreeFlow 项目的实际情况，本文在检索与排序模块中选择了基于元数据的基本检索方式，综合标题、艺人名、描述文本、标签字段、链上身份字段、新鲜度信号与热度信号等因素组织检索结果。这种设计既符合系统当前发展阶段，也有助于后续平滑扩展至更复杂的推荐模型。


\section{研究内容与主要工作}

本文围绕 FreeFlow 链上音乐发行与播放系统开展了相关的研究工作，完成了区块链场景下音乐平台原型的设计与实现，并设计了一套冷启动场景下的音乐检索方案。

本文设计并实现了客户端、链上智能合约、链下Web2.5服务器三部分。客户端使用Electron 和React 构建，为用户完整实现了播放控制、歌单管理、链上控制等模块，并为创作者搭建了额外的交互平台，实现了项目管理、元数据编辑、收益分析等实用界面以方便创作者的交互流程。链上智能合约基于 Solidity 语言和Hardhat 框架开发，实现了用于管理歌曲资源授权信息的MusicAccess1155，用于创作者定义分账方式的分账合约 RoyaltySplitter 以及平台业务方法集合PlatformHub。服务器基于Express 构建，用于实现平台高频业务数据的管理以及链上数据的索引，以实现媲美中心化音乐平台的服务。
功能上，重点关注创作者发布作品和用户购买播放链路，着重介绍这两个场景下的用户体验和界面设计。

本文还设计了一套冷启动场景下的音乐检索方法，以适合冷启动场景下的检索要求，并利用MusicBrainz元数据库进行了评估。







